\chapter{Installation and first run}\label{ch:install}

\section{Install}
RS3D is distributed as source code. It needs Python 3.11 or newer and \href{https://docs.astral.sh/uv/}{\code{uv}}, a tool that
creates an isolated environment with exactly the libraries the application needs.

\begin{minted}{bash}
cd rs3d-app
uv sync
\end{minted}

\code{uv sync} reads the project's lock file and installs the libraries (ObsPy, SciPy, NiceGUI and others) into a folder
named \file{.venv}. It does not touch the rest of your system. It needs the Internet once, the first time.

\section{Start the application}
\begin{minted}{bash}
uv run rs3d-app
\end{minted}

Your browser opens on the application. The first time (no Shake added yet) you land on \ui{Settings}; afterwards you land on
\ui{Status}. To stop the application press \key{Ctrl}+\key{C} in the terminal where you started it.

\begin{rs3dtip}[Choosing the port]
Without options the application listens on port 8765, or on any free port if that one is taken, and prints the address. Use
\code{uv run rs3d-app --port 9000} to choose one, and \code{uv run rs3d-app --no-browser} to start it without opening a window
(for example on a second computer).
\end{rs3dtip}

\section{Where your data is kept}
Everything the application stores lives in one folder, shown on the \ui{Settings} page next to \ui{Data folder}.

\begin{center}\small
\begin{tabular}{@{}ll@{}}
\toprule
System & Default folder\\ \midrule
Linux & \code{\textasciitilde/.local/share/rs3d}\\
macOS & \code{\textasciitilde/Library/Application Support/rs3d}\\
Windows & \code{C:\textbackslash Users\textbackslash <you>\textbackslash AppData\textbackslash Local\textbackslash rs3d}\\
\bottomrule
\end{tabular}
\end{center}

To use a different folder, for example on a larger disk, set the environment variable \code{RS3D\_DATA\_DIR} before starting:

\begin{minted}{bash}
RS3D_DATA_DIR=/mnt/seismic/rs3d uv run rs3d-app
\end{minted}

Chapter~\ref{ch:storage} explains what is inside the folder and how to back it up.

\section{Updating and uninstalling}
To update, replace the program folder with the newer one and run \code{uv sync} again; your data folder is separate and is kept. To
remove the application, delete the program folder and, if you also want to delete your data, the data folder.

\begin{rs3dwarn}[The data folder is your recordings]
Deleting the data folder deletes every day file you downloaded. The Shake itself keeps only the most recent days
(7 by default), so data that exist only on your laptop cannot be fetched again.
\end{rs3dwarn}
